\cotefile{157-5}
\scanpage{558}
\mspage{469}
\unsure{dit|de} $N=\{\varnothing\}$, on voit que la condition supplémentaire
\deleted{\unsure{de sa|de la}} stabilité sur $\unsure{\underline T|\overline{\underline I}|\overline{\underline J}}$ n'est pas
superflue. (\uline{NB} Elle sert à assurer que
$\underline S$ est $L_\pi$-stable dans $M$ pour $\pi$ donné.)

4.16.3. Questions \unsure{apparentées} : j'avais espéré
que pour un foncteur accessible entre
catégories accessibles
\[
f:\unsure{N}\mathrel{\unsure{-|\longrightarrow}}M,
\]
la $\unsure{\leftrightarrows|\rightleftarrows}$-image essentielle, formée des $y\in\operatorname{Ob}M$
qui sont isomorphes à un facteur direct
d'un objet de la forme $f(x)$, où $x\in\operatorname{Ob}N$,
serait accessible dans $N$. L'exemple
ci-dessus (nullement \unsure{éliminatoire} !) de
la sous-catégorie $M_z$ de $M$ (formée des
$x\in\operatorname{Ob}M$ tels que $\operatorname{Hom}_M(x,z)\ne\varnothing$) montre
qu'il n'en est rien. En effet, $M_z$ est
l'image essentielle, et aussi l'image $\leftrightarrows$-essentielle,
du foncteur
\[
M_{/z}\longrightarrow M,
\]
où $M_{/z}$ est accessible et le foncteur précédent
accessible, si $M$ est \unsure{accessible} (\unsure{vérification}
par \unsure{AQT}, cf plus \deleted{\unsure{bas}} énoncé plus général ?).
% archive sheet number: 557
